\chapter{CSR file and trap handling}
\label{ch:csr}

\section{Implemented CSRs}

\file{hdl/csr\_file.v} implements the seven CSRs the brief required, plus
identification registers, plus the counters. Reads are combinational; software
writes commit at the end of S2.

\begin{table}[H]
\centering
\caption{CSR map as implemented. \sig{RO} means a write raises illegal
instruction; \sig{WARL} means a write is accepted and discarded.}
\label{tab:csrmap}
\begin{tabular}{@{}TllL{0.42\linewidth}@{}}
\toprule
\textnormal{\textbf{Addr}} & \textbf{Name} & \textbf{Access} & \textbf{Reset / notes} \\
\midrule
0x300 & \reg{mstatus}   & RW & \texttt{0x0000\_1800}. Live fields: MIE[3], MPIE[7], MPP[12:11] read-only \texttt{11}. \\
0x301 & \reg{misa}      & WARL & \texttt{0x4000\_0100}. MXL = 32-bit, extension I. \\
0x304 & \reg{mie}       & RW & 0. Only bits [31:16] exist; writes to [15:0] are dropped. \\
0x305 & \reg{mtvec}     & RW & 0. BASE[31:2], MODE[1:0]. \\
0x340 & \reg{mscratch}  & RW & 0. Never touched by hardware. \\
0x341 & \reg{mepc}      & RW & 0. Bits [1:0] forced to zero on every write. \\
0x342 & \reg{mcause}    & RW & 0. Bit [31] = interrupt. \\
0x343 & \reg{mtval}     & RW & 0. Fault address, or instruction word on illegal. \\
0x344 & \reg{mip}       & RO & 0. One-hot on [31:16], driven by the PIC offer. \\
0xB00 & \reg{mcycle}    & RW & 0. Lower half of a 64-bit counter. \\
0xB02 & \reg{minstret}  & RW & 0. Lower half; counts committed instructions only. \\
0xB03--0xB07 & \reg{mhpmcounter3..7} & RW & 0. Event counters, lower halves. \\
0xB08--0xB1F & \reg{mhpmcounter8..31} & WARL & 0. Hardwired zero. \\
0xB80 & \reg{mcycleh}   & RW & 0. Upper half of \reg{mcycle}. \\
0xB82 & \reg{minstreth} & RW & 0. Upper half of \reg{minstret}. \\
0xB83--0xB87 & \reg{mhpmcounter3h..7h} & RW & 0. Upper halves. \\
0xB88--0xB9F & \reg{mhpmcounter8h..31h} & WARL & 0. Hardwired zero. \\
0x323--0x33F & \reg{mhpmevent3..31} & WARL & 0. Hardwired zero; events are fixed in hardware. \\
0xF11 & \reg{mvendorid} & RO & 0, an unregistered vendor. \\
0xF12 & \reg{marchid}   & RO & 0. \\
0xF13 & \reg{mimpid}    & RO & 0. \\
0xF14 & \reg{mhartid}   & RO & \sig{HART\_ID} parameter. \\
\bottomrule
\end{tabular}
\end{table}

\subsection{Field detail}

\begin{table}[H]
\centering
\caption{\reg{mstatus} (0x300) and \reg{mtvec} (0x305) fields}
\label{tab:csrfields}
\begin{tabular}{@{}llL{0.58\linewidth}@{}}
\toprule
\textbf{Bits} & \textbf{Field} & \textbf{Behaviour} \\
\midrule
\multicolumn{3}{@{}l}{\itshape mstatus} \\
{[3]}      & MIE  & Global interrupt enable. Cleared by hardware on trap entry. \\
{[7]}      & MPIE & MIE before the trap. \sig{MRET} copies it back into MIE and sets MPIE to 1. \\
{[12:11]}  & MPP  & Read-only \texttt{11}; the core is M-mode only. \\
others   & ---  & Not implemented, read 0. \\
\midrule
\multicolumn{3}{@{}l}{\itshape mtvec} \\
{[31:2]}   & BASE & Handler base, 4-byte aligned. \\
{[1:0]}    & MODE & \texttt{00} direct: every trap goes to BASE. \texttt{01} vectored: interrupts go to
                  \texttt{BASE + 4$\cdot$cause}, exceptions still to BASE. \\
\bottomrule
\end{tabular}
\end{table}

\begin{deviation}
\textbf{\reg{mtvec.MODE} is WARL-constrained on the write path.} MODE is a WARL
field with only two defined encodings, \texttt{00} and \texttt{01}. An earlier
version stored a written \texttt{10} or \texttt{11} verbatim and returned it on a
read. The trap behaviour was direct mode either way, so nothing misbehaved, but
software that probes \reg{mtvec} to discover which modes the core implements was
told it supports one it does not. A reserved MODE is now folded to \texttt{00} at
the write, and the CSR WARL bench checks all four encodings.
\end{deviation}

\subsection{Counters}

All seven counters are architecturally 64-bit and are accessed on RV32 through a
base / base+0x80 pair, as the privileged specification requires. M-mode may write
either half. The update rule is shared by all of them
(\sig{cnt\_next} in \file{hdl/csr\_file.v}):

\begin{center}
\begin{minipage}{0.80\linewidth}
\ttfamily\small
inc = cur + ev;\\
if (wr\_lo)      next = \{ inc[63:32], new\_value(inc[31:0])  \};\\
else if (wr\_hi) next = \{ new\_value(inc[63:32]), inc[31:0] \};\\
else             next = inc;
\end{minipage}
\end{center}

The half not being written always takes the increment, so a carry out of the
written half is still counted. Within the written half the CSR operation
decides. \sig{CSRRW} replaces the value outright and the event of that cycle is
genuinely gone, which is what writing a counter means. \sig{CSRRS} and
\sig{CSRRC} set and clear bits on top of the incremented value, so they keep it.
An earlier version applied all three to the pre-increment value and the
documentation claimed no event was ever lost, which was true only when the write
happened to carry.

\begin{table}[H]
\centering
\caption{Event sources of the performance counters}
\label{tab:hpm}
\begin{tabular}{@{}llL{0.48\linewidth}@{}}
\toprule
\textbf{Counter} & \textbf{Event} & \textbf{What it measures} \\
\midrule
\reg{mcycle}       & every cycle & Elapsed time. \\
\reg{minstret}     & \sig{retire} & Instructions that actually committed. \\
\reg{mhpmcounter3} & \sig{ev\_mispredict} & Mispredict redirects; each costs one cycle. \\
\reg{mhpmcounter4} & \sig{ev\_ibus\_wait} & Cycles S2 was empty waiting for a fetch. \\
\reg{mhpmcounter5} & \sig{ev\_dbus\_stall} & Cycles stalled on a data transaction. \\
\reg{mhpmcounter6} & \sig{trap\_set} & Trap entries of any kind. \\
\reg{mhpmcounter7} & \sig{ev\_wfi\_sleep} & Cycles asleep in \sig{WFI}. \\
\bottomrule
\end{tabular}
\end{table}

The counters that are not provided are tied off rather than left unimplemented:
they read zero and swallow writes instead of trapping, which is what the
specification permits for counters an implementation does not have.
\reg{mcountinhibit} (0x320) is absent by choice --- the specification makes it
optional and defines the absent case as ``all counters run'', which is this
design. The practical consequence is that software cannot freeze a counter around
a measurement window.

\section{Access rules}

Access to a CSR that is not implemented, or an effective write to a read-only
one, raises an illegal instruction rather than a bus error. CSRs live inside the
core and are not on any bus, so a bus error code would be the wrong abstraction,
and silently ignoring the access would hide software defects.

The specification's exception is honoured: \sig{CSRRS} and \sig{CSRRC} with
\sig{rs1} = \sig{x0}, and their immediate forms with a zero immediate, count as
pure reads and do not trap on read-only registers. This is handled upstream ---
\sig{csr\_wen} is never asserted for those forms.

\section{Trap entry and return}

Trap entry and \sig{MRET} are sequenced by \file{cpu\_top} and always win over a
same-cycle software write, so a trapping CSR instruction can never commit its own
write.

\begin{table}[H]
\centering
\caption{Atomic state change on trap entry and return}
\label{tab:trapseq}
\begin{tabular}{@{}lL{0.62\linewidth}@{}}
\toprule
\textbf{Event} & \textbf{State written, all on one clock edge} \\
\midrule
Trap entry & \reg{mepc} $\leftarrow$ return PC, \reg{mcause} $\leftarrow$ cause,
             \reg{mtval} $\leftarrow$ fault payload, \reg{MPIE} $\leftarrow$ \reg{MIE},
             \reg{MIE} $\leftarrow$ 0. \\
\sig{MRET} & \reg{MIE} $\leftarrow$ \reg{MPIE}, \reg{MPIE} $\leftarrow$ 1. \\
\bottomrule
\end{tabular}
\end{table}

There is no cycle in which the state is half updated. In the same cycle the
offending instruction is squashed on its way into DX/WB, so it never reaches
writeback, and fetch is redirected to the trap vector. S3 is untouched --- the
instruction ahead had already committed correctly.

\subsection{Return address}

Exceptions record the offending instruction's PC, so software can inspect, repair
and re-run it. Interrupts record the PC of the instruction to resume: the victim
was cancelled before doing anything, so re-running it is safe. The one special
case is an interrupt that wakes a \sig{WFI}, which records \sig{wfi+4}; otherwise
\sig{MRET} would land back on the \sig{WFI} and sleep again.

\subsection{Trap inside a trap}

Trap entry clears \reg{MIE}, so no interrupt can be taken inside a handler until
\sig{MRET}. Synchronous exceptions inside a handler remain possible, which is
correct --- a defective handler should still fault. At system level the PIC keeps
the serviced source on its nesting stack until the end-of-interrupt, so it cannot
re-offer itself.

\section{Trap-tracking flags}

Trap tracking in \file{cpu\_top.v} is a stack rather than a pair of flags: one
bit per open trap level, recording whether that level was an interrupt, plus a
single signal saying whether any level is open at all.

\begin{table}[H]
\centering
\caption{Trap-tracking flags}
\label{tab:trapflags}
\begin{tabular}{@{}lL{0.66\linewidth}@{}}
\toprule
\textbf{Flag} & \textbf{Purpose} \\
\midrule
\sig{trap\_kind\_q} & One bit per open trap level, 1 if that level was an interrupt.
  Pushed on every accepted trap, popped on every \sig{MRET}. Gates
  \sig{cpu\_irq\_eoi}, so an \sig{MRET} closing a synchronous trap pops no PIC nesting
  level and each nested interrupt handler returns its own end-of-interrupt.
  Depth 16, the largest \reg{NEST\_MAX} the PIC accepts. \\
\sig{cpu\_in\_trap} & High while any trap level is open; it drops on the \sig{MRET}
  that closes the last one, so a nested trap inside a handler keeps it high.
  Exported for observability and assertions only; the PIC does not consume it. \\
\bottomrule
\end{tabular}
\end{table}

\begin{figure}[H]
\centering
\begin{tikzpicture}[x=1mm,y=1mm,>={Stealth[length=2.4mm]},semithick]
  \node[stbi] (run)  at (0,0)   {RUNNING\\\scriptsize depth $=0$};
  \node[stb]  (open) at (92,0)  {TRAP OPEN\\\scriptsize depth $\ge 1$};

  % entry: one transition, because the two trap kinds differ only in the bit
  % that is pushed - not in what the pipeline does
  \draw[->,draw=accent] (run) to[bend left=16]
    node[lbl,pos=0.5,above,align=center]
    {trap taken: push its kind\\
     \scriptsize 1 for an interrupt, 0 for a synchronous trap} (open);

  % the return that closes the last open level
  \draw[->,draw=accent] (open) to[bend left=16]
    node[lbl,pos=0.5,below,align=center]
    {MRET at depth 1: pop\\
     \scriptsize eoi pulses only if the popped bit is 1} (run);

  \draw[->,draw=accent,looseness=6,out=125,in=55] (open) to
    node[lbl,above=0.5mm,align=center]{nested trap: push again} (open);

  \draw[->,draw=accent,looseness=6,out=-55,in=-125] (open) to
    node[lbl,below=0.5mm,align=center]{MRET at depth $>1$: pop, stay open} (open);
\end{tikzpicture}
\caption{CPU trap tracking. The open trap levels are a stack of kind bits, not a
pair of flags: every trap pushes one bit and every \sig{MRET} pops one, so a
synchronous trap taken inside an interrupt handler returns without an
end-of-interrupt, and each nested interrupt handler emits its own.}
\label{fig:trapfsm}
\end{figure}

\section{Cause set}

\begin{table}[H]
\centering
\caption{Supported trap causes and the \reg{mtval} payload, from the priority
chain in \file{hdl/exception\_unit.v}}
\label{tab:causes}
\begin{tabular}{@{}ccL{0.36\linewidth}L{0.12\linewidth}@{}}
\toprule
\textbf{Code} & \textbf{Prio.} & \textbf{Cause and condition} & \textbf{\reg{mtval}} \\
\midrule
0  & 5 & Instruction address misaligned: a taken branch or jump whose target has
       bit~1 set. Bit~0 is zero by construction --- \sig{JALR} masks it and the
       other formats cannot encode it. & target \\
1  & 1 & Instruction access fault: the fetch returned SLVERR or DECERR and the
       word reached S2. & fetch PC \\
2  & 2 & Illegal instruction: unknown encoding, or a disallowed CSR access. & instruction \\
3  & 3 & Breakpoint: \sig{EBREAK}. & PC \\
4  & 6 & Load address misaligned. Checked before the access, so no transaction
       is issued. & address \\
5  & 8 & Load access fault: read response SLVERR or DECERR. & address \\
6  & 7 & Store address misaligned, as cause 4. & address \\
7  & 9 & Store access fault: write response SLVERR or DECERR. & address \\
11 & 4 & Environment call from M-mode: \sig{ECALL}. & 0 \\
16--31 & --- & External interrupt, source 0--15. \reg{mcause}[31] is set and
       source $v$ gives cause $16+v$. Interrupts are decided in
       \file{cpu\_top.v} before this chain runs, not inside it. & 0 \\
\bottomrule
\end{tabular}
\end{table}

The \textbf{Prio.} column is the position in the \texttt{if}/\texttt{else}
chain, which is not the numeric order of the codes. Reading the chain out: a
fetch that failed outranks everything, because the instruction word cannot be
trusted; then illegal, breakpoint and environment call, which are properties of
the decoded instruction; then instruction-address misalignment, which needs the
branch resolved; then the two misalignment causes, which are known before the
access is issued; and last the two access faults, which are only knowable once
the bus has answered. One chain selects the cause and the \reg{mtval} payload
together, so the cause list exists in exactly one place in the source.

The two orderings that matter are checked as tests rather than asserted here:
an address that is both misaligned and unmapped must raise cause 4 and issue
nothing (Section~\ref{sec:vp-traps}), and every cause is given its own isolated
case with all four architectural results compared.

Two rules shape it. Misalignment is checked \emph{before} the access, so a bad
address never produces an AXI transaction at all. Bus errors are checked
\emph{after} the response and become access faults rather than illegal
instructions --- the instruction was perfectly legal, the bus failed.

The three access faults (1, 5, 7) exist because this CPU sits on a real bus that
can refuse a transaction. Without them a bus error would have to be reported as
an illegal instruction, which would mislead the handler.

\section{Vectored mode}

With \reg{mtvec.MODE} = \texttt{01}, interrupts jump to
\texttt{BASE + 4$\cdot$cause}, so PIC source $v$ lands at
\texttt{BASE + 4$\cdot$(16+v)} --- its own slot, normally holding a jump to the
handler. Exceptions still go to BASE. Both paths are exercised separately in the
test program.
